% Running example and continuation.

\centering
\renewcommand{\ULthickness}{0.6pt}
\begin{tikzpicture}[setA, x=1cm, y=1cm]

% Layout coordinates.
\def\xCa{1.50}
\def\xCb{4.00}
\def\xCc{6.50}
\def\xCd{9.50}
\def\xCe{15.05}                   % later frontier c_5 (continuation)
\def\axisY{0}
\def\writeY{0.85}
\def\brackY{-0.30}
\def\brackLabelY{-0.60}
\def\obsY{-1.55}

% Replay/source equality at the frontier.
\def\eqY{3.55}                    % equality unit, center y
\def\eqX{\xCd}                    % equality unit, center x (the frontier)

\node[font=\sffamily\Large, anchor=center] (eq) at (\eqX, \eqY) {$=$};

\node[blockFill, anchor=east, text width=29mm, align=center,
      left=4pt of eq.west, font=\sffamily\footnotesize]
  (lbox) {%
  \textbf{$\Apply(\CleanPrefix) \restrict K$}\\[1pt]
  \begin{tabular}{@{\ }r@{\,:\,}l@{\ }}
    100 & $3000$  \\
    200 & $12000$ \\
    300 & $3500$  \\
  \end{tabular}};

\node[block, anchor=west, text width=29mm, align=center,
      right=4pt of eq.east, font=\sffamily\footnotesize]
  (rbox) {%
  \textbf{$\Src(b_0, H, f) \restrict K$}\\[1pt]
  \begin{tabular}{@{\ }r@{\,:\,}l@{\ }}
    100 & $3000$  \\
    200 & $12000$ \\
    300 & $3500$  \\
  \end{tabular}};

\node[sublabelDim, anchor=south] at ($(lbox.north -| eq) + (0,1.5pt)$)
  {at frontier $f = c_4$, on scope $K = \{100, 200, 300\}$};

% Connect the frontier to the state comparison.
\coordinate (eqSouth) at ($(eq |- lbox.south) + (0, -2pt)$);
\draw[guide, line width=0.5pt]
  (\xCd, \axisY) -- (eqSouth);

% Source-log axis.
\draw[axisLine] (-0.10, \axisY) -- (15.65, \axisY);

\foreach \x/\lab in {\xCa/c_1, \xCb/c_2, \xCc/c_3} {
  \draw[axisTick] (\x, \axisY+0.06) -- (\x, \axisY-0.06);
  \node[anchor=north, font=\sffamily\scriptsize, inner sep=1pt]
    at (\x, \axisY-0.06) {$\lab$};
}
\draw[axisTick, line width=0.8pt] (\xCd, \axisY+0.10) -- (\xCd, \axisY-0.10);
\node[anchor=north, font=\sffamily\scriptsize\bfseries, inner sep=1pt]
  at (\xCd, \axisY-0.08) {$c_4 = f$};

\node[sublabelDim, anchor=west, inner sep=0pt]
  at (15.80, \axisY) {source-log};

% Source writes.
\foreach \x/\op/\key/\val/\tag in {%
  \xCa/Upd/100/$3000$/w_1,
  \xCb/Ins/300/$2500$/w_2,
  \xCc/Upd/200/$12000$/w_3,
  \xCd/Upd/300/$3500$/w_4%
} {
  \node[pillFill, anchor=south] (w-\tag) at (\x, \writeY+0.10) {%
    \strut\textbf{\op}~\key\\\,$\to$~\val\,};
  \draw[guide, line width=0.35pt] (\x, \axisY+0.07) -- (\x, \writeY+0.10);
}

% Chunk reads.
% Chunk A observations.
\draw[braceMirror] (\xCa-0.65, \brackY) -- (\xCa+0.65, \brackY);
\node[sublabel, anchor=north] (chAlab) at (\xCa, \brackLabelY)
  {chunk $A$ at $c_1$};

\node[pill, anchor=north, text width=24mm, align=left, inner xsep=4pt]
  (chAobs) at (\xCa, \obsY) {%
  \,$100 \to 3000$\\
  \,\sout{$200 \to 10000$}};

% Chunk B observations.
\draw[braceMirror] (\xCc-0.65, \brackY) -- (\xCc+0.65, \brackY);
\node[sublabel, anchor=north] (chBlab) at (\xCc, \brackLabelY)
  {chunk $B$ at $c_3$};

\node[pill, anchor=north, text width=24mm, align=left, inner xsep=4pt]
  (chBobs) at (\xCc, \obsY) {%
  \,\sout{$300 \to 2500$}};

% Connect stale observations to the writes that supersede them.
\draw[weakFlow, rounded corners=2pt]
  ($(chAobs.east) + (0.05, -0.10)$)
  to[out=15, in=-100, looseness=0.7]
  ($(w-w_3.south) + (-0.05, 0.02)$);
\draw[weakFlow, rounded corners=2pt]
  ($(chBobs.east) + (0.05, 0)$)
  to[out=15, in=-100, looseness=0.7]
  ($(w-w_4.south) + (-0.05, 0.02)$);

% "dominated by" labels sit just below each superseding write.
\node[font=\sffamily\scriptsize\itshape, text=gMid, anchor=north,
      inner sep=1pt, fill=white]
  at ($(w-w_3.south) + (0, -0.06)$) {dominated by};
\node[font=\sffamily\scriptsize\itshape, text=gMid, anchor=north,
      inner sep=1pt, fill=white]
  at ($(w-w_4.south) + (0, -0.06)$) {dominated by};

% Continuation to the later frontier.

\draw[gLight, line width=0.4pt, densely dashed]
  (\xCe, \axisY) -- (\xCe, \eqY-0.45);

\node[pill, draw=gLight, text=gMid, anchor=south] (w-w5)
  at (\xCe, \writeY+0.10) {\strut\textbf{Del}~200\\\,$\to$~absent\,};
\draw[guide, gLight, line width=0.35pt]
  (\xCe, \axisY+0.07) -- (\xCe, \writeY+0.10);

\draw[axisTick, gMid, line width=0.8pt]
  (\xCe, \axisY+0.10) -- (\xCe, \axisY-0.10);
\node[anchor=north, font=\sffamily\scriptsize\bfseries, text=gMid,
      inner sep=1pt] at (\xCe, \axisY-0.08) {$c_5 = f'$};

\draw[braceMirror, draw=gLight, dashed]
  (\xCd+0.70, \brackY) -- (\xCe-0.70, \brackY);
\node[sublabelDim, anchor=north]
  at ($(\xCd, \brackLabelY)!0.5!(\xCe, \brackLabelY)$)
  {continuation segment: one CDC event};

\node[blockGhost, anchor=center, text width=28mm, align=center,
      font=\sffamily\scriptsize] (c5card) at (\xCe, \eqY) {%
  \textbf{cut advances to $f' = c_5$}\\[1.5pt]
  no new chunk read};

\end{tikzpicture}

\caption{Two chunk reads at different source-log coordinates assemble
into a replay whose state at frontier $f = c_4$ equals the source on
scope $K = \{100, 200, 300\}$. Stale observations (strikethrough) are
dominated by later source-log writes. No single physical read gathers
every key in scope, so the cut is virtual rather than physical.
That is the bootstrap cut at $c_4$. The lighter elements to the right
show the continuation of
Section~\ref{sec:advancing-cut-running-example}. One further
source-log event ($c_5$, a delete of account 200) is appended as the
continuation segment for the half-open interval $(c_4, c_5]$. It begins
strictly after the frontier $c_4$ and requires no additional chunk read. This advances
the same cut to the later frontier $f' = c_5$.}
\label{fig:running-example-timeline}
